% GENERATED FILE. Edit canonical lessons, then run npm run generate:books.
% canonical-source-hash: fnv1a64:7ad42562aff101a1
% canonical-lessons: HI-C321-hamla, HI-C321-gherna, HI-C321-top, HI-C321-bhirna, HI-C321-jujhna

\chapter{The Attack}
\label{ch:hi-the-attack}
\hlchaptermodality{321}

\begin{chapteropening}
\textbf{By the end of this chapter:} \emph{I can tell how an army attacked a fort.}

Complete the last lesson of chapter 321: I can tell how an army attacked a fort.
\end{chapteropening}

\section[\texorpdfstring{\hi{हमला} (hamlā)}{हमला (hamlā)}]{\hi{हमला} (hamlā) — an attack}

\label{lesson:HI-C321-hamla}



\begin{quote}
Before the new one: say the Hindi for a crown, then the Hindi for a royal court.
\end{quote}

\subsection*{You'll want to know: \hi{हमला}}
\textbf{\hi{हमला}} — \emph{hamlā} — \textquotedblleft{}an attack\textquotedblright{}.

\begin{quote}
\hi{दुश्मन} \hi{ने} \hi{रात} \hi{को} \hi{क़िले} \hi{पर} \hi{हमला} \hi{किया।}
\end{quote}

\begin{quote}
\emph{duśman ne rāt ko qile par hamlā kiyā.}
\end{quote}

\begin{quote}
The enemy attacked the fort at night.
\end{quote}

Masculine.

\begin{cousinweb}
\textbf{\hi{हमला}} (\emph{hamlā}) is usually traced, through Persian, to Arabic \emph{ḥamla}, \textquotedblleft{}a charge, an onset\textquotedblright{}. What is attacked takes \textbf{\hi{पर}}: \textbf{\hi{क़िले} \hi{पर} \hi{हमला}}.
\end{cousinweb}

\begin{grammarlens}[title={What you've built}]
One more word for telling what happened.
\end{grammarlens}

\subsection*{Your turn}
\begin{itemize}
\raggedright
  \item \emph{Say it:} \emph{hamlā}
  \item \emph{Say it:} \emph{hamlā}, once more
  \item \emph{Say it:} \emph{duśman ne rāt ko qile par hamlā kiyā.} — The enemy attacked the fort at night.
  \item \emph{Recall:} say the Hindi for a crown, then the Hindi for a royal court, then say \emph{hamlā} again
\end{itemize}

\begin{tcolorbox}[breakable,colback=teal!4,colframe=teal!35!black,title={Before you move on}]
What does \hi{हमला} mean? (\textquotedblleft{}An attack\textquotedblright{}.) Say it once more.
\end{tcolorbox}
\section[\texorpdfstring{\hi{घेरना} (ghernā)}{घेरना (ghernā)}]{\hi{घेरना} (ghernā) — to surround; to lay siege to}

\label{lesson:HI-C321-gherna}



\begin{quote}
Before the new one: say the Hindi for a royal court, then the Hindi for an attack.
\end{quote}

\subsection*{You'll want to know: \hi{घेरना}}
\textbf{\hi{घेरना}} — \emph{ghernā} — \textquotedblleft{}to surround; to lay siege to\textquotedblright{}.

\begin{quote}
\hi{दुश्मन} \hi{की} \hi{सेना} \hi{ने} \hi{क़िले} \hi{को} \hi{घेर} \hi{लिया।}
\end{quote}

\begin{quote}
\emph{duśman kī senā ne qile ko gher liyā.}
\end{quote}

\begin{quote}
The enemy army laid siege to the fort.
\end{quote}

Past, with \hi{ने} (he or she …): \hi{उसने} \hi{गाँव} \hi{को} \hi{घेरा}, \emph{usne gā̃v ko gherā}.

\begin{cousinweb}
\textbf{\hi{घेरना}} (\emph{ghernā}) is inherited from \emph{*ghir}, \textquotedblleft{}go round\textquotedblright{}, a form CDIAL reconstructs (CDIAL 4474): to go all the way round something and close it in.
\end{cousinweb}

\begin{grammarlens}[title={What you've built}]
One more word for telling what happened.
\end{grammarlens}

\subsection*{Your turn}
\begin{itemize}
\raggedright
  \item \emph{Say it:} \emph{ghernā}
  \item \emph{Say it:} \emph{ghernā}, once more
  \item \emph{Say it:} \emph{duśman kī senā ne qile ko gher liyā.} — The enemy army laid siege to the fort.
  \item \emph{Recall:} say the Hindi for a royal court, then the Hindi for an attack, then say \emph{ghernā} again
\end{itemize}

\begin{tcolorbox}[breakable,colback=teal!4,colframe=teal!35!black,title={Before you move on}]
What does \hi{घेरना} mean? (\textquotedblleft{}To surround; to lay siege to\textquotedblright{}.) Say it once more.
\end{tcolorbox}
\section[\texorpdfstring{\hi{तोप} (top)}{तोप (top)}]{\hi{तोप} (top) — a cannon}

\label{lesson:HI-C321-top}



\begin{quote}
Before the new one: say the Hindi for an attack, then the Hindi for to surround.
\end{quote}

\subsection*{You'll want to know: \hi{तोप}}
\textbf{\hi{तोप}} — \emph{top} — \textquotedblleft{}a cannon\textquotedblright{}.

\begin{quote}
\hi{क़िले} \hi{की} \hi{दीवार} \hi{पर} \hi{चार} \hi{तोपें} \hi{थीं।}
\end{quote}

\begin{quote}
\emph{qile kī dīvār par cār topeṁ thīṁ.}
\end{quote}

\begin{quote}
There were four cannons on the fort wall.
\end{quote}

Feminine.

\begin{cousinweb}
\textbf{\hi{तोप}} (\emph{top}) is usually traced, through Persian, to Turkish \emph{top}, \textquotedblleft{}a ball\textquotedblright{}, and so a gun that fires balls. It is feminine in Hindi: \textbf{\hi{तोपें} \hi{थीं}}.
\end{cousinweb}

\begin{grammarlens}[title={What you've built}]
One more word for telling what happened.
\end{grammarlens}

\subsection*{Your turn}
\begin{itemize}
\raggedright
  \item \emph{Say it:} \emph{top}
  \item \emph{Say it:} \emph{top}, once more
  \item \emph{Say it:} \emph{qile kī dīvār par cār topeṁ thīṁ.} — There were four cannons on the fort wall.
  \item \emph{Recall:} say the Hindi for an attack, then the Hindi for to surround, then say \emph{top} again
\end{itemize}

\begin{tcolorbox}[breakable,colback=teal!4,colframe=teal!35!black,title={Before you move on}]
What does \hi{तोप} mean? (\textquotedblleft{}A cannon\textquotedblright{}.) Say it once more.
\end{tcolorbox}
\section[\texorpdfstring{\hi{भिड़ना} (bhiṛnā)}{भिड़ना (bhiṛnā)}]{\hi{भिड़ना} (bhiṛnā) — to clash; to come to blows}

\label{lesson:HI-C321-bhirna}



\begin{quote}
Before the new one: say the Hindi for to surround, then the Hindi for a cannon.
\end{quote}

\subsection*{You'll want to know: \hi{भिड़ना}}
\textbf{\hi{भिड़ना}} — \emph{bhiṛnā} — \textquotedblleft{}to clash; to come to blows\textquotedblright{}.

\begin{quote}
\hi{दोनों} \hi{सेनाएँ} \hi{नदी} \hi{के} \hi{किनारे} \hi{भिड़ीं।}
\end{quote}

\begin{quote}
\emph{donõ senāẽ nadī ke kināre bhiṛīṁ.}
\end{quote}

\begin{quote}
The two armies clashed on the riverbank.
\end{quote}

Past (he …): \hi{वह} \hi{सिपाही} \hi{भिड़} \hi{गया}, \emph{vah sipāhī bhiṛ gayā}.

\begin{cousinweb}
\textbf{\hi{भिड़ना}} (\emph{bhiṛnā}) is inherited from \emph{*bhiṭ}, a form CDIAL reconstructs for meeting and crowding together (CDIAL 9490): two sides that \textbf{\hi{भिड़ते} \hi{हैं}} come right up against each other.
\end{cousinweb}

\begin{grammarlens}[title={What you've built}]
One more word for telling what happened.
\end{grammarlens}

\subsection*{Your turn}
\begin{itemize}
\raggedright
  \item \emph{Say it:} \emph{bhiṛnā}
  \item \emph{Say it:} \emph{bhiṛnā}, once more
  \item \emph{Say it:} \emph{donõ senāẽ nadī ke kināre bhiṛīṁ.} — The two armies clashed on the riverbank.
  \item \emph{Recall:} say the Hindi for to surround, then the Hindi for a cannon, then say \emph{bhiṛnā} again
\end{itemize}

\begin{tcolorbox}[breakable,colback=teal!4,colframe=teal!35!black,title={Before you move on}]
What does \hi{भिड़ना} mean? (\textquotedblleft{}To clash; to come to blows\textquotedblright{}.) Say it once more.
\end{tcolorbox}
\section[\texorpdfstring{\hi{जूझना} (jūjhnā)}{जूझना (jūjhnā)}]{\hi{जूझना} (jūjhnā) — to fight hard; to struggle}

\label{lesson:HI-C321-jujhna}



\begin{quote}
Before the new one: say the Hindi for a cannon, then the Hindi for to clash.
\end{quote}

\subsection*{You'll want to know: \hi{जूझना}}
\textbf{\hi{जूझना}} — \emph{jūjhnā} — \textquotedblleft{}to fight hard; to struggle\textquotedblright{}.

\begin{quote}
\hi{सिपाही} \hi{तीन} \hi{दिन} \hi{तक} \hi{दुश्मन} \hi{से} \hi{जूझे।}
\end{quote}

\begin{quote}
\emph{sipāhī tīn din tak duśman se jūjhe.}
\end{quote}

\begin{quote}
The soldiers fought the enemy for three days.
\end{quote}

Past (he …): \hi{वह} \hi{सिपाही} \hi{जूझा}, \emph{vah sipāhī jūjhā}.

\begin{cousinweb}
\textbf{\hi{जूझना}} (\emph{jūjhnā}) goes back to Sanskrit \emph{yudhyate}, \textquotedblleft{}fights\textquotedblright{} (CDIAL 10502), the family of \textbf{\hi{युद्ध}}, war. It is a long, hard fight, against an enemy or against trouble.
\end{cousinweb}

\begin{grammarlens}[title={What you've built}]
One more word for telling what happened.
\end{grammarlens}

\subsection*{Your turn}
\begin{itemize}
\raggedright
  \item \emph{Say it:} \emph{jūjhnā}
  \item \emph{Say it:} \emph{jūjhnā}, once more
  \item \emph{Say it:} \emph{sipāhī tīn din tak duśman se jūjhe.} — The soldiers fought the enemy for three days.
  \item \emph{Recall:} say the Hindi for a cannon, then the Hindi for to clash, then say \emph{jūjhnā} again
\end{itemize}

\begin{tcolorbox}[breakable,colback=teal!4,colframe=teal!35!black,title={Before you move on}]
What does \hi{जूझना} mean? (\textquotedblleft{}To fight hard; to struggle\textquotedblright{}.) Say it once more.
\end{tcolorbox}
